\documentclass[10pt,a4paper]{amsart}
\usepackage[margin=19mm]{geometry}
\usepackage{amsmath,amssymb,mathtools}
\usepackage[scr=rsfs,cal=euler]{mathalfa}
\usepackage{xfrac}
\usepackage[unicode,hidelinks,bookmarks=false]{hyperref}
\newcommand{\ie}{i.e.\ }
\newcommand{\cf}{cf.\ }
\newcommand{\resp}{respectively\ }
\newcommand{\Subsection}{\subsection}
\providecommand{\nobreakdash}{\nobreak\hbox{-}}
\setlength{\emergencystretch}{2em}
\multlinegap0pt
\usepackage{xcolor}
\usepackage{amssymb}
\newtheorem{thm}{Theorem}
\def\nn{\nonumber}
\newcommand{\oL }{\mathcal{L}(E) }
\title{Matrix-valued truncated Toeplitz operators: Properties and Applications}
\author{Rewayat Khan, Flavia Colonna}
\date{Source: arXiv:2609.15480v1 (CC BY 4.0)}
\begin{document}
\maketitle

\noindent\emph{Source and licence.} Matrix-valued truncated Toeplitz operators: Properties and Applications. Version arXiv:2609.15480v1 (CC BY 4.0), distributed by arXiv under the Creative Commons Attribution 4.0 International licence, http://creativecommons.org/licenses/by/4.0/, as recorded in the arXiv licence metadata for this exact version and shown on the abstract page https://arxiv.org/abs/2609.15480v1. Full licence notice: \texttt{licenses/arxiv-2609.15480-CC-BY-4.0.txt}.\par\smallskip

\noindent\textit{Editorial scope.} Counted unit: the article's thm G (source0.tex lines 953--960). The article's complete original proof of it is delivered with it (source0.tex lines 961--1050, one \texttt{proof} environment of 90 lines). No mathematics is written, repaired or re-derived: the statement and the proof are the article's own lines, in the article's own notation, and no retained line is substituted or re-flowed.  No further source-local context is needed: every cross-reference of every retained line resolves inside the two delivered spans.  Nothing was omitted from any retained span, so the package carries no omission record.  No retained line contains a citation, so the wrapper carries no bibliography and no bibliography entry.  No retained line contains a figure environment, an image-inclusion command or drawing code, so no image is delivered and the package declares no asset dependency.  Editorial formatting only, in addition: an AMS \texttt{amsart} preamble that restates the article's own declarations and macros used by the retained lines, namely the environments \texttt{thm} on separate counters, together with the standard packages those lines need.  The retained environments print in this document as Theorem 1 in order of appearance, whereas the article prints the same items with the numbering of its own full document.  The complete native source of the article as distributed in the arXiv e-print for this exact version is bundled unchanged as \texttt{sources/210418/source0.tex}.
\begin{thm}\label{G} Let $\Theta_1$ and $\Theta_2$ be inner functions and 
 $\Phi\in L^{\infty}(\oL)$. Then $A_{\Phi}^{\Theta_{1},\Theta_{2}}$ is invertible if and only if the Toeplitz operator $T_{G}: H^2(E)\oplus H^2(E)\to H^2(E)\oplus H^2(E)$ is invertible,
where $G=\begin{pmatrix}
\Theta_{2}& \Phi\\
0&\Theta_{1}^{*}
\end{pmatrix}.
$
\end{thm}

\begin{proof}
Consider the equation generated by the operator $A_{\Phi}^{\Theta_{1},\Theta_{2}}\in \mathcal{T}(\Theta_{1},\Theta_{2})$, i.e.,
\begin{equation}\label{asymm}
A_{\Phi}^{\Theta_{1},\Theta_{2}}f_{\Theta_{1}}=P_{\Theta_{2}}(\Phi f_{\Theta_{1}})=g_{\Theta_{2}},
\end{equation}
where $f_{\Theta_{1}}\in K_{\Theta_{1}}$ and $g_{\Theta_{2}}\in K_{\Theta_{2}}$.
\medskip

The equation generated by Toeplitz operator $T_{G}$ with block matrix symbol is given by
 \begin{equation}\label{TG}
P_{+}\begin{pmatrix}
\Theta_{2}& \Phi\\
0&\Theta_{1}^{*}
\end{pmatrix}
\begin{pmatrix}
X_{1}\\
X_{2}
\end{pmatrix}
=\begin{pmatrix}
f_{1}\\
f_{2}
\end{pmatrix},
\end{equation}
where $X_{1}, X_{2}, f_{1}, f_{2}\in H^{2}(E)$.

Suppose $A_{\Phi}^{\Theta_{1},\Theta_{2}}:K_{\Theta_{1}}\to K_{\Theta_{2}}$ is invertible. We wish to prove that \eqref{TG} has unique solution. Performing block multiplication, this equation can be written as the system of two equations
%
 %We can write  \eqref{TG} as follows:
 \begin{equation}\label{first}
 P_{+}(\Theta_{2}X_{1})+P_{+}(\Phi X_{2})=f_{1}
 \end{equation}
 \begin{equation}\label{second}
 P_{+}(\Theta^{*}_{1}X_{2})=f_{2}.
 \end{equation}
 The general solution of \eqref{second} has the  form
 $$X_{2}=\psi+\Theta_{1}f_{2},$$
 where $\psi$ is an arbitrary function in %from
  $K_{\Theta_{1}}$. Substituting this into %Then
  \eqref{first} yields %has the form
 \begin{equation}\label{third}
 \Theta_{2}X_{1}+P_{+}(\Phi\psi)+P_{+}(\Phi\Theta_{1}f_{2})=f_{1}.
 \end{equation}
 Applying the projection operator $P_{\Theta_{2}}$ to both sides of \eqref{third}, we obtain
  \begin{equation}%\label{4th}
 P_{\Theta_{2}}P_{+}(\Phi\psi)+P_{\Theta_{2}}P_{+}(\Phi\Theta_{1}f_{2})=P_{\Theta_{2}}(f_{1}).\nn
 \end{equation}
Equivalently % It can be written as
 \begin{equation}\label{5th}
 A_{\Phi}^{\Theta_{1},\Theta_{2}} \psi=P_{\Theta_{2}}(f_{1}-P_{+}(\Phi\Theta_{1}f_{2})).
 \end{equation}
 By the invertibility of $A_{\Phi}^{\Theta_{1},\Theta_{2}}$,  %implies that
  equation \eqref{5th} has unique solution $\psi\in K_{\Theta_{1}}$, which in turn determines $X_2$. Since, $T_{\Theta_{2}}$ is isometric with range equal to $\Theta_2 H^2(E)$, by \eqref{third}, we have
  $$
 T_{\Theta_{2}}X_{1}=f_1-P_{+}(\Phi\psi)-P_{+}(\Phi\Theta_{1}f_{2})=h\in \Theta_2 H^2(E).
 $$
The equation
 \begin{equation}\label{X_1h}
     T_{\Theta_2}X_1=h
 \end{equation}
has at most one solution. 
Since $T_{\Theta_2}$  is injective and $h$ is in its range, there is unique $X_1\in H^2(E)$ satisfying \eqref{X_1h}. Explicitly, $$X_1=\Theta_2^{*} h\in H^2(E)$$
satisfying
 $$T_{\Theta_2}(X_1)=T_{\Theta_2}(\Theta_2^{*}h)=h.$$ 
 
Conversely, assume \eqref{TG} has unique solution for arbitrary $f_{1}, f_{2}\in H^{2}(E)$. We may express  \eqref{asymm} in the form
 \begin{equation}\label{atto}
 P_{+}(\Phi f_{\Theta_{1}})+\Theta_{2}X_{1}=g_{\Theta_{2}},
 \end{equation}
 where $X_{1}$ is some unknown function in %from
  $H^{2}(E)$. Along with \eqref{atto}, we consider an equation of the kind
 \begin{equation}\label{zero}
 P_{+}(\Theta_{1}^{*}f_{\Theta_{1}})=0.
 \end{equation}
 Rewriting \eqref{atto} and \eqref{zero} of the form:
 \begin{equation}\label{last}
P_{+}\begin{pmatrix}
\Theta_{2}& \Phi\\
0&\Theta_{1}^{*}
\end{pmatrix}
\begin{pmatrix}
X_{1}\\
f_{\Theta_{1}}
\end{pmatrix}
=\begin{pmatrix}
g_{\Theta_{2}}\\
0
\end{pmatrix},
\end{equation}
such that \eqref{last} has unique solution, we deduce that the operator $A_{\Phi}^{\Theta_{1},\Theta_{2}}$ is invertible. Hence, the required result follows.
\end{proof}

\end{document}
